\section{Зачем нужны большие числа}

Наивная мечта «придумать самое большое число» на поверку оказывается вполне
математической. Приведём аргументы, выгодно выставляющие эту детскую идею.

\begin{enumerate}
  \item \textbf{Давать имя необозримому --- это и есть математика.} Гугол, гуголплекс,
  число Грэма, \(\mathrm{TREE}(3)\) --- математика постоянно присваивает компактные имена
  и символы конечным величинам, которые никто не в силах выписать. Придумать символ для
  огромного числа --- рефлекс математика, а не каприз ребёнка.
  \item \textbf{Дети переоткрыли тетрацию.} Правило \(L^{L}\) --- это тетрация, а лестница
  \(a \dots z\) --- быстрорастущая иерархия, родственная функции Аккермана, стрелкам
  Кнута и быстрорастущим иерархиям из теории доказательств.
  \item \textbf{Именованная лестница --- опора для интуиции.} Размеченная ступенями
  \(a \dots z\), она даёт чувство масштаба, недоступное одному чудовищному числу.
  \item \textbf{Это входная дверь.} Игра «а какое число самое большое?» на лавочке во
  дворе для многих --- первый настоящий вкус математики. Ей и посвящена эта работа.
\end{enumerate}

\begin{table}[h]
\centering
\caption{Несколько реальных больших чисел и их порядок \(\log_{10}\), для сравнения.}
\begin{tabular}{lrl}
\toprule
Величина & \(\log_{10}\) & Где встречается \\
\midrule
Ундециллион (\(=a\)) & \startlog & начало нашей лестницы \\
Число атомов во Вселенной & \(\approx 80\) & космология \\
Число Шеннона (шахматы) & \(\approx 120\) & теория игр \\
Гугол & \(100\) & популяризация \\
Планковских объёмов во Вселенной & \(\approx 185\) & физика \\
Уровень \(b\) & \(\bmagnitude\) & \emph{за пределами всего физического} \\
\bottomrule
\end{tabular}
\label{tab:scales}
\end{table}

Уже уровень \(b\) оставляет позади любую физическую величину. Наша лестница --- не про
подсчёт чего-либо реального, а про саму идею восхождения к бесконечности.
